\songtitle{Tu puro vivir}

\tag*{2026}\tag{salsa}\tag{spanish-lyrics}
\begin{notes}
Reimagining of the Anthem of Bogotá as a salsa song.

Original lyrics by Pedro Medina Avendaño, with minor alterations by Carlos Thompson for genre adaptation.
\end{notes}
\begin{style}
salsa dura, 122 BPM, piano tumbao, claves, timbales, cowbell, bass montuno, gang choir, male tenor lead, female soprano lead, brass attack, call and response, spoken intro, 1970s live mix, plate reverb, analog tape saturation, crowd stomp chorus, minor key, triumphant energy
\end{style}
\begin{lyrics}
\songpart{Intro}
Tres guerreros abrieron tus ojos\\
a una espada, a una cruz y a un pendón.\\
Desde entonces no hay miedo en tus lindes\\
ni codicia en tu gran corazón.

\songpart{Verse 1}
Hirió el hondo diamante un agosto\\
el cordaje de un nuevo laúd\\
y hoy se escucha el fluir melodioso\\
en los himnos de la juventud.

Fértil madre de altiva progenie\\
que sonríe ante el vano oropel\\
siempre atenta a la luz de la mañana\\
y al pasado y su luz siempre fiel.

\songpart{Chorus}
Entonemos un himno a tu cielo\\
a tu tierra y tu puro vivir\\
¡Blanca estrella que alumbra en los Andes!\\
¡Ancha senda que va al porvenir!
(¡Ancha senda que va al porvenir!)

\songpart{Verse 2}
La sabana es un cielo caído\\
una alfombra tendida a tus pies\\
y del mundo variado que animas\\
eres brazo y cerebro a la vez.

Sobreviven de un reino dorado\\
de un imperio sin puestas de sol\\
en ti un templo, un escudo, una reja\\
un retablo, una pila, un farol.

\songpart{Chorus}
Entonemos un himno a tu cielo\\
a tu tierra y tu puro vivir\\
¡Blanca estrella que alumbra en los Andes!\\
¡Ancha senda que va al porvenir!
(¡Ancha senda que va al porvenir!)

\songpart{Verse 3}
Al gran Caldas que escruta los astros\\
y a Bolívar que torna a nacer\\
a Nariño accionando la imprenta\\
como en sueños los vuelves a ver.

Caros, Cuervos y Pombos y Silvas\\
tantos nombres de fama inmortal\\
que en el hilo sin fin de la historia\\
les dio vida tu amor maternal.

\songpart{Chorus}
Entonemos un himno a tu cielo\\
a tu tierra y tu puro vivir\\
¡Blanca estrella que alumbra en los Andes!\\
¡Ancha senda que va al porvenir!
(¡Ancha senda que va al porvenir!)

\songpart{Bridge}
Oriflama de la Gran Colombia\\
en Caracas y Quito estará\\
para siempre en la luz de tu gloria\\
con las dianas de la libertad.

Noble, leal en la paz y en la guerra\\
de tus fueres colinas al pie\\
y en el arco de la media lunar\\
resucitas en Cid, Santa Fe.

\songpart{Chorus}
Entonemos un himno a tu cielo\\
a tu tierra y tu puro vivir\\
¡Blanca estrella que alumbra en los Andes!\\
¡Ancha senda que va al porvenir!

\songpart{Outro}
Flor de razas compendio y corona\\
en la patria no hay otra ni habrá\\
nuestra voz la repiten los siglos:\\
¡Bogotá! ¡Bogotá! ¡Bogotá!
\end{lyrics}

